% !TeX root = ../../../金融经济学_王江_讲义及习题.tex
% 左图沿用两基金分离示意图的双曲线和效用形式，坐标按版面缩放。
\begin{tikzpicture}[
 >=Stealth,x=1.25cm,y=1.12cm,font=\small,
 frontier/.style={blue!70!black,very thick},
 guide/.style={gray,densely dotted},
 choice/.style={teal!65!black},
 second/.style={red!68!black}]
 \draw[->] (0,0)--(3.55,0) node[right] {$\sigma$};
 \draw[->] (0,0)--(0,2.9) node[above] {$\bar r$};
 \coordinate (A) at ({1.34*sqrt(1+.48^2)},{.72*(.9-.48)});
 \coordinate (Pone) at ({1.34*sqrt(1+.65^2)},{.72*(.9+.65)});
 \coordinate (Ptwo) at ({1.34*sqrt(1+1.75^2)},{.72*(.9+1.75)});
 \coordinate (B) at ({1.34*sqrt(1+2.08^2)},{.72*(.9+2.08)});
 \draw[guide] (0,.3024) node[left,black] {$\bar r_A$}--(A);
 \draw[guide] (0,2.1456) node[left,black] {$\bar r_B$}--(B);
 % 复用章末两基金图的四次效用曲线，按本图坐标缩放。
 \pgfmathdeclarefunction{wjindifference}{2}{%
  \pgfmathparse{.72*(.9+#2+((#1/1.34)^4-(1+#2^2)^2)/(4*#2*(1+#2^2)))}}
 \foreach \offset in {-.28,-.14,.14,.28} {
  \draw[teal!45!white,dashed] plot[domain=.84:1.96,samples=41,variable=\s]
   (\s,{wjindifference(\s,.65)+\offset});
  \draw[red!45!white,dashed] plot[domain=1.83:3.12,samples=41,variable=\s]
   (\s,{wjindifference(\s,1.75)+\offset});
 }
 \draw[choice,thick] plot[domain=.84:1.96,samples=41,variable=\s]
  (\s,{wjindifference(\s,.65)});
 \draw[second,thick] plot[domain=1.83:3.12,samples=41,variable=\s]
  (\s,{wjindifference(\s,1.75)});
 \draw[frontier] plot[domain=-.7:2.34,samples=81,variable=\t]
  ({1.34*sqrt(1+\t*\t)},{.72*(.9+\t)});
 \fill[blue!70!black] (A) circle (2pt) node[right=3pt,black] {$A$};
 \fill[blue!70!black] (B) circle (2pt) node[above right=3pt,black] {$B$};
 \fill[choice] (Pone) circle (2pt) node[above left=2pt,fill=white,inner sep=1pt] {$P_1$};
 \fill[second] (Ptwo) circle (2pt) node[below left=2pt,fill=white,inner sep=1pt] {$P_2$};
 \node at (1.68,-.72) {(a) 证券期望收益率不全相同};

 \begin{scope}[xshift=6.05cm]
  \draw[->] (0,0)--(3.55,0) node[right] {$\sigma$};
  \draw[->] (0,0)--(0,2.9) node[above] {$\bar r$};
  \draw[guide] (0,1.32) node[left,black] {$a/c$}--(.78,1.32)--(.78,0)
   node[below,black] {$1/\sqrt c$};
  % 两人的最高可达无差异曲线都经过 MVP，斜率不同。
  \foreach \offset in {-.28,-.14,.14,.28} {
   \draw[teal!45!white,dashed] plot[domain=.30:2.30,samples=41,variable=\s]
    (\s,{1.32+.25*(\s^2-.78^2)+\offset});
   \draw[red!45!white,dashed] plot[domain=.30:1.60,samples=41,variable=\s]
    (\s,{1.32+.65*(\s^2-.78^2)+\offset});
  }
  \draw[choice,thick] plot[domain=.30:2.30,samples=61,variable=\s]
   (\s,{1.32+.25*(\s^2-.78^2)});
  \draw[second,thick] plot[domain=.30:1.60,samples=61,variable=\s]
   (\s,{1.32+.65*(\s^2-.78^2)});
  \draw[gray!80,thick,<-] (.78,1.32)--(3.3,1.32)
   node[pos=.70,below=4pt,black] {同均值的可行组合};
  \fill[black] (.78,1.32) circle (2.3pt);
  \node[align=center] (common) at (2.12,.48)
    {$P_1=P_2=\mathrm{MVP}$};
  \draw[gray!80,->] (common.north west)--(.83,1.27);
  \node at (1.68,-.72) {(b) 证券期望收益率相同};
 \end{scope}
\end{tikzpicture}
